\chapter{Introduction}
\thispagestyle{plain}
\label{ch:intro}

\section{Background}
Small unmanned aerial vehicles (UAVs) are now used in groups, or swarms, for
surveillance, search and rescue, disaster assessment and border patrol. Each UAV
carries a camera and collects images of the area it covers. Training across the
datasets of several UAVs may improve generalization compared with training on a
single UAV's data \cite{mcmahan2017fedavg}, but sending raw imagery to a ground
station is costly in bandwidth and exposes sensitive data if a link is intercepted.

Federated learning (FL) addresses this by moving training to the data. Each UAV
trains a copy of the model on its own images and uploads only the resulting model
update. A server (here, a ground base station) combines the updates into a new
global model and sends it back. The raw images remain on the UAV, and an update is
usually much smaller than the data it was computed from. Updates can still leak some
information about the data, so FL alone is not a privacy guarantee.

A tactical swarm, however, breaks two assumptions that standard FL relies on.
\begin{enumerate}
\item \textbf{Links are unreliable.} Air-to-ground channels fade as the UAV moves,
  hand over between base stations, and may be deliberately jammed. Parts of an update
  are lost, and some updates miss the round deadline entirely.
\item \textbf{Some clients may be hostile.} A captured or spoofed UAV can follow the
  protocol while sending a poisoned update designed to damage the global model. Such
  clients are called \emph{Byzantine}.
\end{enumerate}
Most Byzantine-robust aggregation rules assume that every update arrives intact, and
most work on FL over imperfect channels assumes that every client is honest. A swarm
needs both at once.

\section{Problem Statement}
Given a UAV swarm that trains a shared image classifier with hierarchical FL over
lossy air-to-ground links, in which a minority of UAVs are Byzantine, design and
evaluate an aggregation rule that (i) keeps the global model accurate and stable,
(ii) does not penalize honest UAVs simply because their links are poor, and
(iii) can be compared fairly against established rules under identical conditions.

\section{Aim and Scope}
The aim of this project is to determine \emph{which} kind of channel awareness
actually helps a robust aggregator in a UAV swarm, and to quantify the effect with a
controlled, reproducible experiment. Its scope in Phase~I is:
\begin{itemize}
\item a simulation of hierarchical FL with up to three base stations, UAV
  association and handover;
\item packet-level uplink losses from a synthetic jammed channel and from a measured
  UAV received-power trace;
\item five supported attacks (label flip, sign flip, Gaussian noise, ALIE and IPM),
  of which the preliminary study uses label flip and the controlled study the three
  model-poisoning attacks sign flip, ALIE and IPM;
\item six aggregation rules evaluated on the same data, models, attackers and channel
  realizations, and the base paper's two erasure rules for a separate comparison.
\end{itemize}
Hardware flight tests, real radio transmission and attacks that forge link telemetry
are outside the Phase~I scope.

\section{Objectives}
\begin{enumerate}
\item Build an FL baseline for a UAV swarm on aerial-transformed CIFAR-10 and measure its
  accuracy under IID and non-IID data.
\item Model a lossy uplink in which lost packets erase parts of an update, and add
  Byzantine attackers.
\item Design a reputation-based aggregator (CCRA) that uses a trusted reference
  update, and give it two forms of channel awareness: erasure masking and
  telemetry-weighted reputation.
\item Separate the effect of each component with a prespecified, multi-seed ablation
  study, and report the results with honest uncertainty estimates.
\end{enumerate}

\section{Contributions}
\begin{itemize}
\item An open, tested PyTorch framework for hierarchical UAV-swarm FL with
  trace-driven packet erasures, handover and five attacks, with automated tests.
\item A preliminary study, then a controlled \NumRuns-run study whose analysis plan
  and decision threshold were fixed before the runs.
\item The finding that erasure masking improves accuracy under attack on weak links
  by \MaskGain{} percentage points, while telemetry-weighted reputation lowers it,
  together with an analysis of the reputation rule that suggests a mechanism
  consistent with this result.
\item A completed \BpNumRuns-run comparison of the base paper's two erasure rules
  on matched seeds, attackers and channels, with validated checkpoint recovery and
  an explicit separation from the prespecified ablation.
\item A four-page manuscript in IEEE letter format that reports these results.
\end{itemize}

\section{Organization of the Report}
Chapter~\ref{ch:lit} reviews federated learning, Byzantine-robust aggregation and FL
over unreliable wireless links, and identifies the gap this project addresses.
Chapter~\ref{ch:method} presents the theoretical study: the system and threat
models, the aggregation rules and an analysis of the reputation update.
Chapter~\ref{ch:results} presents the experimental study: implementation, protocol,
the baseline, the preliminary study and the controlled study, and a discussion of
what they show. Chapter~\ref{ch:conclusions} states the conclusions and the plan for
Phase~II. The appendices, after the references, list the experimental settings and
the steps needed to reproduce every result.
